\documentclass[11pt,a4paper]{article}
\usepackage[margin=1in]{geometry}
\usepackage{xcolor}
\usepackage{listings}
\usepackage{tikz}
\usepackage{hyperref}
\title{Vortex Compiler (VXC)\\Language Specification and Phase 1 Lexer}
\author{Vikash Kumar Gupta}
\date{\today}
\definecolor{codegray}{gray}{0.95}
\lstset{basicstyle=\ttfamily\small,backgroundcolor=\color{codegray},frame=single,breaklines=true}
\begin{document}
\maketitle
\section{Purpose}
Vortex Compiler (VXC) is an educational compiler implemented in C. Its initial target is a custom stack-based virtual machine. Phase 1 defines the language and implements lexical analysis.
\section{Compilation Pipeline}
\begin{center}
\begin{tikzpicture}[node distance=8mm, every node/.style={draw,rounded corners,align=center,minimum width=4cm,minimum height=8mm}]
\node (s) {Source Code}; \node (l) [below=of s] {Lexer}; \node (p) [below=of l] {Parser / AST};\node (sem) [below=of p] {Semantic Analysis};\node (ir) [below=of sem] {Intermediate Representation};\node (opt) [below=of ir] {Optimization};\node (vm) [below=of opt] {Virtual Machine};
\foreach \a/\b in {s/l,l/p,p/sem,sem/ir,ir/opt,opt/vm} {\draw[->] (\a)--(\b);}
\end{tikzpicture}
\end{center}
\section{Lexical Specification}
The lexer recognizes whitespace, line comments beginning with \texttt{//}, identifiers, decimal integer literals, keywords, and punctuation.
\subsection{Keywords}
\texttt{fn}, \texttt{let}, \texttt{int}, and \texttt{print} are reserved keywords.
\subsection{Token Grammar}
\begin{lstlisting}
identifier   := letter (letter | digit | '_')*
integer      := digit+
letter       := 'A'..'Z' | 'a'..'z' | '_'
digit        := '0'..'9'
\end{lstlisting}
\section{Example Program}
\begin{lstlisting}
fn main() {
    let x: int = 10 + 20;
    print(x);
}
\end{lstlisting}
\section{Phase 1 Implementation}
The current C implementation reads a source file and prints each token with its type, lexeme, line, and column. This provides a testable foundation for the parser.
\section{Build and Run}
\begin{lstlisting}
make
./build/vxc examples/hello.vxc
\end{lstlisting}
\section{Next Phase}
Phase 2 will add a recursive-descent parser and an abstract syntax tree for declarations, expressions, and function bodies.
\end{document}
